\chapter{System Design and Implementation}
\label{ch:design}

\section{Functional Design and Sub-problem Partitioning}

The project was partitioned by evidence and operating layer rather than by one
monolithic model:

\begin{enumerate}
  \item \textbf{Capture and preprocessing:} record a short clip, sample frames,
        locate/crop/normalise faces, and reject unusable capture.
  \item \textbf{Mobile temporal screening:} compute per-frame MobileNetV3 features,
        aggregate across the clip, and return a first-layer probability.
  \item \textbf{Rich temporal analysis:} combine visual tokens and facial-landmark
        motion in $G^2V^2$former.
  \item \textbf{Spatial domain generalisation:} evaluate GD-FAS under source-only
        and fine-tuned transfer protocols.
  \item \textbf{Deepfake specialisation:} train a universal Xception baseline and
        one binary specialist per manipulation method, then learn a meta-router.
  \item \textbf{Policy and governance:} calibrate scores, select escalation bands,
        audit decisions, preserve privacy, and provide human fallback.
\end{enumerate}

Temporal information appears at multiple levels, illustrated in
\Cref{fig:temporal-levels}. The lightweight model averages frame embeddings; the
transformer represents explicit time/landmark relationships; and the target
decision layer can evaluate score stability across the session. Only the first
two have experimental artefacts in the repository.

\begin{figure}[H]
  \centering
  \resizebox{0.95\textwidth}{!}{\begin{tikzpicture}[
  font=\small,
  box/.style={draw=buetblue,rounded corners,fill=softblue,align=left,text width=118mm,minimum height=12mm},
  arr/.style={-{Latex[length=2.2mm]},thick,draw=buetblue}
]
\node[box] (a) {\textbf{MobileNetV3 temporal aggregation:} per-frame visual features are mean-pooled across a sampled clip before binary classification.};
\node[box,below=7mm of a] (b) {\textbf{$G^2V^2$former:} face appearance and 68-landmark motion are fused through spatial and temporal attention for video-level presentation-attack detection.};
\node[box,below=7mm of b] (c) {\textbf{Decision-level extension:} aggregate calibrated evidence across frames/models and escalate unstable or uncertain cases. This stage is proposed, not integrated.};
\draw[arr] (a) -- (b);
\draw[arr] (b) -- (c);
\end{tikzpicture}

}
  \caption{Temporal evidence investigated or proposed at three levels.}
  \label{fig:temporal-levels}
\end{figure}

\section{Design Refinement and Simulation}

The design evolved through successive findings. Early work treated a small
mobile model as the primary detector and a server as a stronger fallback. Strong
within-domain physical-PAD results motivated explicit temporal research. The
$G^2V^2$former experiments then exposed poor transfer to high-quality digital
manipulations, while GD-FAS provided a spatial domain-generalisation comparison.
Fine-tuning improved some deepfake results but introduced source-retention and
unseen-domain trade-offs. The later Xception study therefore tested whether the
heterogeneity of manipulation artefacts was better handled with specialists.

Two refinements follow from this progression. First, ``liveness'' must include
capture provenance: a visually live face can still be injected. Second, the
second layer must not be described as one fixed model before common-protocol
calibration and ablation show which evidence is complementary. The architecture
therefore presents $G^2V^2$former, GD-FAS, and Xception routing as evaluated
candidates, not an already fused product.

\subsection{Mobile temporal model}

For frame $x_t$, the MobileNetV3 backbone produces feature vector $h_t$. The clip
representation and probability are
\begin{equation}
  \bar h=\frac{1}{K}\sum_{t=1}^{K}h_t,
  \qquad p_1=\sigma(w^\top \bar h+b).
\end{equation}
The exported and application contract uses $K=24$. A training script in the
repository uses $K=10$, so checkpoint provenance and sequence-length consistency
must be reconciled before the exported artefact is associated with a reported
training run. Mean pooling is efficient and permutation-invariant, but it cannot
represent temporal order.

\subsection{\texorpdfstring{$G^2V^2$former}{G2V2Former} and GD-FAS}

$G^2V^2$former uses an eight-frame input in the inspected experiment and 68
facial landmarks per frame. Its explicit spatiotemporal attention makes it a
candidate server model, but face/landmark preprocessing dominates a part of
latency and introduces an additional failure surface. GD-FAS operates on frames
and optimises group-aware/domain-invariant features. The comparison is deliberate:
it asks whether temporal dynamics or spatial domain regularisation transfers more
reliably across the selected datasets.

\subsection{Embedding-conditioned specialist routing}

\begin{figure}[H]
  \centering
  \resizebox{0.98\textwidth}{!}{\begin{tikzpicture}[
  font=\footnotesize,
  node distance=7mm and 8mm,
  box/.style={draw=buetblue,rounded corners,fill=softblue,align=center,minimum height=11mm,text width=25mm},
  expert/.style={draw=gray!80,rounded corners,fill=white,align=center,minimum height=7mm,text width=25mm},
  arr/.style={-{Latex[length=2mm]},thick,draw=buetblue}
]
\node[box] (x) {preprocessed\\face frame $x$};
\node[box,right=of x] (phi) {embedding\\$z=\phi(x)$};
\node[box,right=of phi] (router) {meta-router\\$q_e(z)$};
\node[expert,above right=-1mm and 13mm of router] (e1) {Deepfakes expert};
\node[expert,below=2mm of e1] (e2) {Face2Face expert};
\node[expert,below=2mm of e2] (e3) {FaceShifter expert};
\node[expert,below=2mm of e3] (e4) {FaceSwap expert};
\node[expert,below=2mm of e4] (e5) {NeuralTextures expert};
\node[box,text width=31mm,right=13mm of e3] (out) {$\widehat p(\mathrm{fake}\mid x)$\\from selected expert};

\draw[arr] (x) -- (phi);
\draw[arr] (phi) -- (router);
\foreach \e in {e1,e2,e3,e4,e5}{\draw[arr] (router.east) -- (\e.west);}
\foreach \e in {e1,e2,e3,e4,e5}{\draw[arr] (\e.east) -- (out.west);}
\node[below=6mm of e5,align=center,text width=62mm,text=draftred] {No ``original'' route: every frame is assigned to one binary specialist.};
\end{tikzpicture}
}
  \caption{Intended hard specialist-routing formulation. A meta-model selects one
  of five binary experts from the frame embedding; it never predicts an original
  class.}
  \label{fig:intended-router}
\end{figure}

The specialists share an Xception architecture (approximately 20.81 million
parameters) but are tuned separately for Deepfakes, Face2Face, FaceShifter,
FaceSwap, and NeuralTextures. Xception's depthwise-separable convolution provides
a consistent baseline architecture \cite{chollet2017xception}. The intended hard
routing design in \Cref{fig:intended-router} computes only the selected expert at
decision time if router features are available independently. In the current
notebooks, however, router labels are six-way method/original labels and the
original branch uses $1-P(\mathrm{original})$. That implementation tests a useful
method-conditioned routing hypothesis but not the final mathematical definition.

To train a true best-expert router, predictions for every training frame must be
generated out-of-fold. Let $\ell(s_e(x_i),y_i)$ be expert $e$'s binary loss; then
one possible target is
\begin{equation}
  e_i^*=\arg\min_{e\in\mathcal E}\ell(s_e^{(-f(i))}(x_i),y_i),
\end{equation}
where $s_e^{(-f(i))}$ excludes sample $i$'s fold from its training. A deterministic
tie rule and a policy for examples on which all experts fail are required.

\section{Application of Tools}

\begin{table}[H]
\centering
\caption{Main tools and their engineering roles.}
\label{tab:tools}
\begin{tabularx}{\textwidth}{@{}p{32mm}Y Y@{}}
\toprule
Tool & Role & Use in this project \\
\midrule
PyTorch / torchvision & Training, transfer learning, evaluation, export &
MobileNetV3, $G^2V^2$former, GD-FAS, and Xception experiments \cite{pytorch}. \\
Flutter / Dart & Cross-platform mobile application & Camera capture, 24-frame
preprocessing, PyTorch Lite invocation, and local result display \cite{flutter}. \\
FastAPI & Research inference service & Separately demonstrated endpoint for the
$G^2V^2$former workflow; not connected to Flutter. \\
XGBoost & Router-head ablation & Tree-based method prediction from saved
embeddings in the legacy six-way routing experiment. \\
OpenCV / face and landmark tooling & Frame extraction and facial preprocessing &
Face crops and landmark sequence construction; failure handling needs hardening. \\
Jupyter / Colab & Reproducible experiment narrative & Training, evaluation,
result tables, and latency pilot; environment/version locking remains incomplete. \\
Git / GitHub & Source history and collaboration & Version control for source and
reports; large datasets/checkpoints require separate governed storage. \\
\bottomrule
\end{tabularx}
\end{table}

\section{Implementation}

The repository contains four main implementation strands. The mobile pipeline
includes temporal MobileNetV3 training/export code and a Flutter application. The
$G^2V^2$former strand contains data preparation, adapted training/evaluation, and
a FastAPI prototype. The GD-FAS strand contains domain-generalisation evaluation
and fine-tuning work. The \texttt{DeepFake-MoE} directory contains two Xception
notebooks: the main CNN routing experiment and the XGBoost router-head variant.
The directory name is retained as a repository label; this report uses
``embedding-conditioned specialist routing'' because the current hard, single
expert selection is more precise than mixture-of-experts terminology.

The routing dataset was constructed from FaceForensics++ C23 with original and
five fake methods. The saved configuration uses 1,000 video groups: 900 training
groups and 100 validation groups. Sampling 100 frames per method/video yields
540,000 training frames and 60,000 validation frames. Groups, rather than frames,
are separated to reduce video-content leakage. Models take $299\times299$ images;
training includes blur, JPEG compression, horizontal flipping, AdamW with initial
learning rate $10^{-4}$ and weight decay $10^{-5}$, batch size 128, up to 15
epochs, and a binary threshold of 0.5.

\section{Deployment}

The only integrated deployment artefact is the local Flutter first stage. It
captures a clip, creates one 24-frame input, invokes the PyTorch Lite model, and
returns a thresholded score. The inspected application has server escalation
disabled. The $G^2V^2$former FastAPI service demonstrates that a stronger model
can be exposed independently, not that the two layers form one working system.

A production deployment would require versioned model contracts; authenticated,
encrypted client--server communication; genuine-camera and device-integrity
signals; replay protection; strict media retention; rate limiting; monitoring;
calibrated policies; rollback; and a reviewer/appeal process. Raw face videos
should not be logged by default. Audit records should store the minimum necessary
metadata---such as consent/purpose, model version, quality state, calibrated
scores, decision reason, and override---under a declared retention schedule.
